% !TEX root = ../../Main.tex
\section{Evaluation Protocol}
\label{Section:EvaluationProtocol}

The account, the cost model and the execution engine are described in \Cref{Section:TradingFramework} and are not repeated here. Three points are specific to the evaluation.

The evaluation window runs from 1 January 2015 to 1 January 2026, which is eleven years. It starts one year after the data begins. This way the slowest indicator in the observation, a moving average over five thousand hourly bars, is fully warmed up before the first decision.

Every result is measured from five different starting balances: $9\,900$, $10\,000$, $10\,050$, $10\,100$ and $10\,200$ euros. A strategy whose result depends on its starting point is not a real strategy. Small changes to the opening balance change the rounding of every position size, and so they change the whole sequence of trades. Reporting one balance would hide this sensitivity. The tables therefore give the mean across the five balances and the dispersion around it.

Position size is a parameter set for each pair separately. It was chosen by the sweep described in \Cref{Section:ExperimentalProcedure} and was not fixed in advance. Five pairs use the reference risk of \Cref{Equation:PositionSizing}. NZD/USD uses one and a half times that risk, and USD/JPY uses a smaller fraction of it. Position size scales every trade, but it does not change the policy that chooses the trades. The ratios, the alpha and beta decomposition and the correlation structure are therefore not affected by it.